\documentclass[12pt,a4paper]{article}

% ============================================================================
% MUSTERLÖSUNG: IT-SICHERHEIT KLAUSUR (90 Minuten)
% ============================================================================
% Generiert: 27.05.2026
% Basierend auf Modulhandbuch, Skripten und Notizen
% ============================================================================

\usepackage[utf8]{inputenc}
\usepackage[german]{babel}
\usepackage[a4paper, margin=2cm]{geometry}
\usepackage{fancyhdr}
\usepackage{amsmath}
\usepackage{amssymb}
\usepackage{listings}
\usepackage{xcolor}
\usepackage{graphicx}
\usepackage{tikz}
\usepackage{tabularx}
\usepackage{enumerate}
\usepackage{tcolorbox}

% Code-Styling
\lstset{
    language=Python,
    basicstyle=\ttfamily\small,
    keywordstyle=\color{blue},
    commentstyle=\color{gray},
    stringstyle=\color{red},
    breaklines=true,
    showstringspaces=false
}

% Farben für Lösungen
\definecolor{solutionbg}{gray}{0.95}
\definecolor{pointcolor}{RGB}{0,100,0}

% Header/Footer
\pagestyle{fancy}
\fancyhf{}
\rhead{IT-Sicherheit Musterlösung}
\lhead{DHBW Stuttgart}
\cfoot{\thepage}

% ============================================================================
\title{\textbf{MUSTERLÖSUNG: IT-Sicherheit Klausur}}
\author{DHBW Stuttgart -- Informatik}
\date{27.05.2026}

\begin{document}

\maketitle

% ============================================================================
% HINWEISE FÜR KORREKTOR
% ============================================================================

\section*{HINWEISE FÜR KORREKTOR}

\begin{tcolorbox}[colback=yellow!10, colframe=orange!50, title=Bewertungshinweise]
\begin{itemize}
    \item Die folgende Musterlösung zeigt exemplarische Lösungen für alle Aufgaben
    \item Alternative korrekte Antworten sind zu akzeptieren, wenn diese sachlich richtig und gut begründet sind
    \item Teilpunkte sind großzügig zu vergeben bei Verständnis der Konzepte
    \item Rechenfehler sollten nur dann abgezogen werden, wenn der Rechenweg richtig ist
    \item Kreuze \checkmark markieren Lösungen, die vollständig sind und die volle Punktzahl erhalten
\end{itemize}
\end{tcolorbox}

\newpage

\tableofcontents

\newpage

% ============================================================================
% AUFGABE 1: GRUNDLAGEN
% ============================================================================

\section{Grundlagen der IT-Sicherheit [15 Punkte]}

\subsection{Aufgabe 1.1 - Schutzziele (6 Punkte)}

\subsubsection*{Aufgabenstellung}
Nennen Sie die sechs Schutzziele der IT-Sicherheit und erläutern Sie kurz, was unter dem Schutzziel \textit{Vertraulichkeit} zu verstehen ist. Geben Sie ein Beispiel, wie dieses Schutzziel in der Praxis verletzt werden kann.

\subsubsection*{Musterlösung}

\begin{tcolorbox}[colback=solutionbg, colframe=black!30]

\textbf{Die sechs Schutzziele der IT-Sicherheit:} \checkmark

\begin{enumerate}
    \item \textbf{Vertraulichkeit (Confidentiality)}: Nur autorisierte Personen dürfen Zugriff auf Informationen haben
    \item \textbf{Integrität (Integrity)}: Daten dürfen nicht unbefugt verändert werden
    \item \textbf{Verfügbarkeit (Availability)}: Systeme und Daten müssen verfügbar sein
    \item \textbf{Authentizität (Authenticity)}: Identität von Absender/Empfänger kann überprüft werden
    \item \textbf{Verbindlichkeit/Zurechenbarkeit (Accountability)}: Handlungen können einer Person zugeordnet werden
    \item \textbf{Privatheit/Anonymität (Privacy/Anonymity)}: Daten können nicht auf Personen zurückgeführt werden
\end{enumerate}

\textbf{Erklärung Vertraulichkeit:} \checkmark

Vertraulichkeit bedeutet, dass nur diejenigen Personen Zugriff auf bestimmte Informationen haben dürfen, die dazu berechtigt sind. Unautorisierte Personen sollen diese Informationen nicht einsehen können.

\textbf{Praktisches Beispiel für Verletzung:} \checkmark

Ein Angreifer hackt sich in einen Mailserver ein und liest die Mails der Benutzer, ohne dazu berechtigt zu sein. Dies verletzt die Vertraulichkeit, da private Kommunikation mitgelesen wird.

\textbf{Weiteres Beispiel:}
\begin{itemize}
    \item Man-in-the-Middle Angriff: Abhören unverschlüsselter Kommunikation
    \item Datenleck: Unbefugtes Herunterladen von Kundendaten aus einer Datenbank
    \item USB-Stick mit sensiblen Daten wird in öffentlichen Verkehrsmitteln vergessen
\end{itemize}

\end{tcolorbox}

\textbf{Punkteverteilung:}
\begin{itemize}
    \item Alle 6 Schutzziele genannt (3 Punkte) \checkmark
    \item Vertraulichkeit erläutert (2 Punkte) \checkmark
    \item Praktisches Beispiel (1 Punkt) \checkmark
\end{itemize}

\subsection{Aufgabe 1.2 - Bedrohungsklassifikation (9 Punkte)}

\subsubsection*{Aufgabenstellung}

\begin{enumerate}[(a)]
    \item Klassifizieren Sie folgende Bedrohungen und geben Sie jeweils ein Beispiel an:
    \begin{itemize}
        \item Höhere Gewalt
        \item Fahrlässigkeit
        \item Vorsatz
    \end{itemize}
    
    \item Unterscheiden Sie zwischen \textit{passiven} und \textit{aktiven} Angriffen und nennen Sie je ein Beispiel.
\end{enumerate}

\subsubsection*{Musterlösung}

\begin{tcolorbox}[colback=solutionbg, colframe=black!30]

\textbf{(a) Bedrohungsklassifikation:}

\textbf{1. Höhere Gewalt:} \checkmark
\begin{itemize}
    \item \textbf{Definition:} Natürliche oder außergewöhnliche Ereignisse, die nicht kontrollierbar sind
    \item \textbf{Beispiele:}
    \begin{itemize}
        \item Blitzschlag, der Stromausfall verursacht
        \item Erdbeben, das Rechenzentrum beschädigt
        \item Hochwasser/Überschwemmung
        \item Feuer durch äußere Einwirkung
    \end{itemize}
\end{itemize}

\textbf{2. Fahrlässigkeit:} \checkmark
\begin{itemize}
    \item \textbf{Definition:} Unbeabsichtigte Fehler durch menschliches Versagen
    \item \textbf{Beispiele:}
    \begin{itemize}
        \item Datenbankadministrator löscht versehentlich wichtige Daten
        \item Mitarbeiter gibt Passwort an der Rezeption preis
        \item Unverschlüsseltes Backup wird in öffentliche Cloud hochgeladen
        \item Falsche Firewall-Regel öffnet Port unbeabsichtigt
    \end{itemize}
\end{itemize}

\textbf{3. Vorsatz:} \checkmark
\begin{itemize}
    \item \textbf{Definition:} Absichtliche Angriffe durch böswillige Akteure
    \item \textbf{Beispiele:}
    \begin{itemize}
        \item Hacker führt SQL-Injection durch um Datenbank zu manipulieren
        \item Malware-Infection durch Drive-by-Download
        \item Mitarbeiter stiehlt absichtlich Kundendaten
        \item DDoS-Angriff gegen Website
        \item Ransomware-Attacke auf Unternehmensnetzwerk
    \end{itemize}
\end{itemize}

\end{tcolorbox}
\begin{tcolorbox}[colback=solutionbg, colframe=black!30]

\textbf{(b) Passive vs. Aktive Angriffe:} \checkmark

\textbf{Passive Angriffe:}
\begin{itemize}
    \item \textbf{Definition:} Angreifer versucht Informationen zu gewinnen, ohne System zu verändern
    \item \textbf{Schutzziel:} Vertraulichkeit verletzt
    \item \textbf{Beispiel:} Sniffing-Angriff
    \begin{itemize}
        \item Angreifer setzt seinen Computer in den Promiscuous-Modus
        \item Erfasst alle Netzwerk-Pakete im selben Subnetz
        \item Liest Passwörter aus unverschlüsselter Kommunikation
    \end{itemize}
    \item \textbf{Gegenmassnahmen:} Kryptographie (verschlüsselte Übertragung)
\end{itemize}

\textbf{Aktive Angriffe:}
\begin{itemize}
    \item \textbf{Definition:} Angreifer verändert, erzeugt oder löscht Daten
    \item \textbf{Schutzziele:} Integrität und/oder Verfügbarkeit verletzt
    \item \textbf{Beispiel:} Man-in-the-Middle (MITM) Angriff
    \begin{itemize}
        \item Angreifer positioniert sich zwischen Sender und Empfänger
        \item Fängt Datenpakete ab und verändert diese
        \item Sendet manipulierte Pakete weiter
        \item Empfänger erhält verfälschte Informationen
    \end{itemize}
    \item \textbf{Gegenmassnahmen:} 
    \begin{itemize}
        \item Authentifikation (Absenderidentität überprüfen)
        \item Digitale Signaturen (Tamper-Detection)
        \item Integrität-Prüfung (Checksummen, MACs)
    \end{itemize}
\end{itemize}

\end{tcolorbox}

\textbf{Punkteverteilung:}
\begin{itemize}
    \item Höhere Gewalt klassifiziert + Beispiel (2 Punkte) \checkmark
    \item Fahrlässigkeit klassifiziert + Beispiel (2 Punkte) \checkmark
    \item Vorsatz klassifiziert + Beispiel (2 Punkte) \checkmark
    \item Passive Angriffe erklärt + Beispiel (2 Punkte) \checkmark
    \item Aktive Angriffe erklärt + Beispiel (2 Punkte) \checkmark
\end{itemize}

Allerdings waren nur 9 Punkte zu vergeben, nicht 10. Die Gewichtung sollte sein:
\begin{itemize}
    \item (a) 6 Punkte: Je 2 Punkte pro Klassifikation
    \item (b) 3 Punkte: 1,5 Punkte pro Angrifftyp
\end{itemize}

% ============================================================================

\newpage

\section{Symmetrische Verschlüsselung [20 Punkte]}

\subsection{Aufgabe 2.1 - Caesar-Chiffre (8 Punkte)}

\subsubsection*{Aufgabenstellung}

Gegeben ist eine Caesar-Chiffre mit Schlüssel $k = 3$ und Blockgröße 1 (ECB-Modus).
Verwenden Sie die Zuordnung: A $\to$ 0, B $\to$ 1, ..., Z $\to$ 25

\begin{enumerate}[(a)]
    \item Verschlüsseln Sie die Nachricht ``SICHERHEIT''
    \item Welche Sicherheitsprobleme hat die ECB-Modus Verschlüsselung?
    \item Nennen Sie einen sichereren Betriebsmodus und begründen Sie die Verbesserung
\end{enumerate}

\subsubsection*{Musterlösung}

\begin{tcolorbox}[colback=solutionbg, colframe=black!30]

\textbf{(a) Verschlüsselung von ``SICHERHEIT'' mit Caesar-Chiffre ($k=3$):} \checkmark

Caesar-Chiffre: $c_i = (m_i + k) \bmod 26$

\centering
\begin{tabular}{|c|c|c|c|c|}
    \hline
    \textbf{Klartext} & \textbf{Wert} & \textbf{+3} & \textbf{mod 26} & \textbf{Geheimtext} \\
    \hline
    S & 18 & 21 & 21 & V \\
    I & 8 & 11 & 11 & L \\
    C & 2 & 5 & 5 & F \\
    H & 7 & 10 & 10 & K \\
    E & 4 & 7 & 7 & H \\
    R & 17 & 20 & 20 & U \\
    H & 7 & 10 & 10 & K \\
    E & 4 & 7 & 7 & H \\
    I & 8 & 11 & 11 & L \\
    T & 19 & 22 & 22 & W \\
    \hline
\end{tabular}

\textbf{Geheimtext:} $\boxed{\text{VLFKHUHKLW}}$

\textbf{(b) Sicherheitsprobleme der ECB-Verschlüsselung:} \checkmark

ECB (Electronic Code Book) hat folgende Sicherheitsprobleme:

\begin{enumerate}[1.]
    \item \textbf{Pattern Preservation:}
    \begin{itemize}
        \item Gleiche Klartextblöcke ergeben gleiche Geheimtextblöcke
        \item Im obigen Beispiel: ``H'' (Wert 7) wird immer zu ``K`` (Wert 10) verschlüsselt
        \item Siehe Tabelle: E $\to$ H (2x), H $\to$ K (2x), I $\to$ L (2x)
        \item Ein Angreifer kann Muster erkennen und Informationen ableiten
    \end{itemize}
    
    \item \textbf{Statistische Angriffe:}
    \begin{itemize}
        \item Durch Häufigkeitsanalyse können Schlüssel erraten werden
        \item Buchstaben mit hoher Häufigkeit in Klartexten erzeugen Muster im Geheimtext
        \item Im Englischen: ``E'' ist häufigster Buchstabe $\to$ wird immer gleich verschlüsselt
    \end{itemize}
    
    \item \textbf{Visuelle Erkennbarkeit:}
    \begin{itemize}
        \item Bei Bildern oder Strukturdaten: Muster bleiben erkennbar
        \item Beispiel: Rechtecke in ECB-verschlüsselten Bildern bleiben als Rechtecke erkennbar
    \end{itemize}
\end{enumerate}

\end{tcolorbox}
\begin{tcolorbox}[colback=solutionbg, colframe=black!30]

\textbf{(c) Sicherer Betriebsmodus - CBC (Cipher Block Chaining):} \checkmark

\textbf{Funktionsweise CBC:}
$$c_i = E(m_i \oplus c_{i-1})$$

wobei $c_0 = IV$ (Initialization Vector)

\textbf{Vorteile gegenüber ECB:}

\begin{enumerate}[1.]
    \item \textbf{Keine Pattern-Wiederholung:}
    \begin{itemize}
        \item Gleiche Klartextblöcke erzeugen unterschiedliche Geheimtextblöcke
        \item Durch Abhängigkeit vom vorherigen Block
        \item Auch bei zufälligem IV
    \end{itemize}
    
    \item \textbf{Statistisches Sicherheit:}
    \begin{itemize}
        \item Häufigkeitsanalyse ist nicht möglich
        \item Geheimtext sieht zufällig aus
    \end{itemize}
    
    \item \textbf{Diffusion:}
    \begin{itemize}
        \item Eine Änderung im Klartext beeinflusst alle nachfolgenden Ciphertextblöcke
        \item Tamper-Detection wird möglich
    \end{itemize}
\end{enumerate}

\textbf{Andere sichere Modi:}
\begin{itemize}
    \item \textbf{CTR (Counter Mode):} Blockchiffre als Stromchiffre
    \item \textbf{GCM (Galois/Counter Mode):} Mit authentifiierter Verschlüsselung
    \item \textbf{OFB (Output Feedback):} Ähnlich wie CBC
\end{itemize}

\end{tcolorbox}

\textbf{Punkteverteilung:}
\begin{itemize}
    \item Verschlüsselung korrekt (3 Punkte) \checkmark
    \item ECB-Sicherheitsprobleme (3 Punkte) \checkmark
    \item CBC-Modus + Verbesserung (2 Punkte) \checkmark
\end{itemize}

\subsection{Aufgabe 2.2 - AES und DES (12 Punkte)}

\subsubsection*{Musterlösung}

\begin{tcolorbox}[colback=solutionbg, colframe=black!30]

\textbf{(a) Vergleich DES vs. AES:} \checkmark

\centering
\begin{tabular}{|l|l|l|}
    \hline
    \textbf{Kriterium} & \textbf{DES} & \textbf{AES} \\
    \hline
    \textbf{Schlüssellänge} & 56 Bit & 128, 192, 256 Bit \\
    \hline
    \textbf{Blockgröße} & 64 Bit & 128 Bit \\
    \hline
    \textbf{Anzahl Runden} & 16 & 10, 12, 14 (je nach Schlüssellänge) \\
    \hline
    \textbf{Sicherheit heute} & UNSICHER & SICHER \\
    \hline
    \textbf{Verabschiedung} & 1977 & 2001 (NIST Standard) \\
    \hline
    \textbf{Grundkonzept} & Feistel-Netzwerk & Substitution-Permutation \\
    \hline
\end{tabular}

\textbf{(b) Warum ist DES heute nicht mehr sicher?} \checkmark

\begin{enumerate}[1.]
    \item \textbf{Zu kurze Schlüssellänge (56 Bit):}
    \begin{itemize}
        \item Nur $2^{56}$ mögliche Schlüssel $\approx 7,2 \times 10^{16}$
        \item Moderne Computer können dies in wenigen Stunden durchsuchen
        \item Brute-Force Attacke ist praktisch durchführbar
    \end{itemize}
    
    \item \textbf{Triple-DES (3DES) wurde zur Übergangslösung:}
    \begin{itemize}
        \item Verwendet effektiv 112 Bit Schlüssel
        \item Ist aber deutlich langsamer
        \item Blockgröße bleibt nur 64 Bit (zu klein)
    \end{itemize}
    
    \item \textbf{Kleine Blockgröße (64 Bit):}
    \begin{itemize}
        \item Nur $2^{64}$ verschiedene mögliche Blöcke
        \item Nach ca. $2^{32}$ verschlüsselten Blöcken tritt Birthday-Paradox auf
        \item Wiederholte Blöcke werden sichtbar
    \end{itemize}
    
    \item \textbf{Theoretische Schwächen:}
    \begin{itemize}
        \item Lineare Kryptanalyse ist möglich
        \item Differentielle Kryptanalyse hat Ansatzpunkte
    \end{itemize}
\end{enumerate}

\end{tcolorbox}
\begin{tcolorbox}[colback=solutionbg, colframe=black!30]

\textbf{(c) Empfohlene Standards für Industrie 4.0 und IoT:} \checkmark
    
\begin{enumerate}[\textbf{1.}]
    \item \textbf{AES-256-GCM} (ideal)
    \begin{itemize}
        \item 256-Bit Schlüssel für Zukunftssicherheit
        \item GCM (Galois/Counter Mode) für authentifizierte Verschlüsselung
        \item Schützt vor Tampering und Abhören
        \item Performance ist gut
    \end{itemize}
    
    \item \textbf{ChaCha20-Poly1305}
    \begin{itemize}
        \item Schneller auf Geräten ohne AES Hardware-Beschleunigung
        \item Moderner Stream-Cipher
        \item Verwendet in modernen Protokollen (TLS 1.3)
    \end{itemize}
    
    \item \textbf{Lightweight Kryptographie für Resource-Constrained Devices:}
    \begin{itemize}
        \item Simon und Speck (NSA, aber kontrovers)
        \item PRESENT, RECTANGLE
        \item Weniger Ressourcen benötigt, aber geringere Sicherheit
    \end{itemize}
    
    \item \textbf{Zu vermeiden:}
    \begin{itemize}
        \item DES und 3DES (veraltet)
        \item RC4 (Stromchiffre mit bekannten Schwächen)
        \item WEP (Wireless Encryption Protocol - komplett unsicher)
    \end{itemize}
\end{enumerate}

\end{tcolorbox}

\textbf{Punkteverteilung:}
\begin{itemize}
    \item DES vs. AES Vergleich (4 Punkte) \checkmark
    \item Warum DES unsicher (5 Punkte) \checkmark
    \item Empfehlungen für Industrie 4.0/IoT (3 Punkte) \checkmark
\end{itemize}

% ============================================================================

\newpage

\section{Kryptologie und Hashfunktionen [18 Punkte]}

\subsection{Aufgabe 3.1 - Hashfunktionen (9 Punkte)}

\subsubsection*{Musterlösung}

\begin{tcolorbox}[colback=solutionbg, colframe=black!30]

\textbf{(a) Anforderungen an sichere Hashfunktionen:} \checkmark

Mindestens 3 von folgenden (mit je 1 Punkt):

\begin{enumerate}[1.]
    \item \textbf{Preimage-Resistance (Erste Urbild-Resistenz):}
    \begin{itemize}
        \item Gegeben Hash $h = H(m)$, sollte es praktisch unmöglich sein, $m$ zu finden
        \item Brute-Force würde $2^{n}$ Versuche benötigen (mit $n$ = Hash-Länge)
    \end{itemize}
    
    \item \textbf{Second Preimage-Resistance (Zweite Urbild-Resistenz):}
    \begin{itemize}
        \item Gegeben $m$ und $h = H(m)$, sollte es unmöglich sein, ein anderes $m' \neq m$ mit $H(m') = h$ zu finden
        \item Verhindert Kollisionen für ein bekanntes Eingabewort
    \end{itemize}
    
    \item \textbf{Collision-Resistance (Kollisionsresistenz):}
    \begin{itemize}
        \item Es sollte praktisch unmöglich sein, zwei unterschiedliche Eingaben $m \neq m'$ zu finden mit $H(m) = H(m')$
        \item Brute-Force würde ca. $2^{n/2}$ Versuche benötigen (Birthday-Paradox)
        \item Dies ist die wichtigste Eigenschaft
    \end{itemize}
    
    \item \textbf{Determinismus:}
    \begin{itemize}
        \item Gleiche Eingabe muss immer gleichen Hash ergeben
        \item Notwendig für Verifikation
    \end{itemize}
    
    \item \textbf{Avalanche-Effekt:}
    \begin{itemize}
        \item Kleine Änderung in der Eingabe sollte großen Unterschied im Hash verursachen
        \item Verhindert, dass ähnliche Eingaben ähnliche Hashes haben
    \end{itemize}
    
    \item \textbf{Effizienz:}
    \begin{itemize}
        \item Hash sollte schnell berechnet werden können
        \item Wichtig für praktische Anwendungen
    \end{itemize}
\end{enumerate}

\end{tcolorbox}
\begin{tcolorbox}[colback=solutionbg, colframe=black!30]

\textbf{(b) Angriffe bei fehlender Preimage-Resistance:} \checkmark

Wenn $H$ ist nicht Preimage-resistant:

\begin{enumerate}[1.]
    \item \textbf{Passwort-Recovery:}
    \begin{itemize}
        \item Wenn Passwort-Hash $h = H(pw)$ bekannt ist
        \item Und man kann Preimage finden: $pw = H^{-1}(h)$
        \item Passwort-Datenbank ist kompromittiert
    \end{itemize}
    
    \item \textbf{Signatur-Fälschung:}
    \begin{itemize}
        \item Digitale Signatur schützt nur den Hash, nicht die Nachricht selbst
        \item Wenn man zu einem gegebenen Hash $h$ eine neue Nachricht $m'$ finden kann
        \item Die alte Signatur wird gültig für neue Nachricht
    \end{itemize}
    
    \item \textbf{Integrität-Verletzung:}
    \begin{itemize}
        \item Message Authentication Code (MAC) wird unsicher
        \item Angreifer kann Nachricht ändern und neuen MAC berechnen
    \end{itemize}
\end{enumerate}

\textbf{(c) SHA-1 vs. SHA-256:} \checkmark

\centering
\begin{tabular}{|l|l|l|}
    \hline
    \textbf{Eigenschaft} & \textbf{SHA-1} & \textbf{SHA-256} \\
    \hline
    \textbf{Hash-Länge} & 160 Bit & 256 Bit \\
    \hline
    \textbf{Theoretische Sicherheit} & $2^{80}$ & $2^{128}$ \\
    \hline
    \textbf{Birthday-Attacke} & $2^{80}$ & $2^{128}$ \\
    \hline
    \textbf{Praktische Attacken} & Existieren! & Noch nicht praktisch \\
    \hline
    \textbf{Status} & GEBROCHEN & SICHER \\
    \hline
\end{tabular}

\textbf{Warum SHA-1 unsicher ist:}

\begin{enumerate}[1.]
    \item \textbf{Zu kurze Hash-Länge:}
    \begin{itemize}
        \item 160 Bit $\approx 2^{80}$ Möglichkeiten
        \item Moderne Computer können $2^{60}$ Operationen pro Jahr durchführen
        \item Nach Birthday-Paradox: ca. $2^{80}$ Hashes genügen für Kollision
        \item Realistische Attacke in Supercomputer-Zeit
    \end{itemize}
    
    \item \textbf{Schwächen in der Struktur:}
    \begin{itemize}
        \item SHAttered-Angriff (2017): Google demonstrierte erste praktische Kollision
        \item Verwendet SHA1('X.509')+SHA1('PEM')
        \item Bewies, dass Collision-Resistance verloren gegangen ist
    \end{itemize}
    
    \item \textbf{Zukunftssicherheit:}
    \begin{itemize}
        \item Mit 256 Bit ist auch nach 20 Jahren noch keine praktische Attacke absehbar
        \item Quantum-Computer sind noch nicht vorhanden
        \item SHA-256 ist Post-Quantum sicher genug (vorerst)
    \end{itemize}
\end{enumerate}

\end{tcolorbox}

\textbf{Punkteverteilung:}
\begin{itemize}
    \item Mindestens 3 Anforderungen (3 Punkte) \checkmark
    \item Angriffe (3 Punkte) \checkmark
    \item SHA-1 vs SHA-256 Vergleich (3 Punkte) \checkmark
\end{itemize}

\subsection{Aufgabe 3.2 - Message Authentication Code (9 Punkte)}

\subsubsection*{Musterlösung}

\begin{tcolorbox}[colback=solutionbg, colframe=black!30]

\textbf{(a) HMAC vs. einfache Keyed-Hash-Verfahren:} \checkmark

\textbf{Einfache Keyed-Hash Varianten:}

\begin{enumerate}[1.]
    \item \textbf{Secret-Prefix:} $MAC(m, k) = H(k \| m)$
    \item \textbf{Secret-Suffix:} $MAC(m, k) = H(m \| k)$
    \item \textbf{Secret-Sandwich:} $MAC(m, k) = H(k \| m \| k)$
\end{enumerate}

\textbf{HMAC-Struktur:}
$$HMAC(m, k) = H((k \oplus opad) \| H((k \oplus ipad) \| m))$$

wobei:
\begin{itemize}
    \item $opad$ = äußeres Padding (0x5C bytes)
    \item $ipad$ = inneres Padding (0x36 bytes)
\end{itemize}

\textbf{Unterschiede:} \checkmark

\centering
\begin{tabular}{|l|l|l|}
    \hline
    \textbf{Aspekt} & \textbf{Einfache Verfahren} & \textbf{HMAC} \\
    \hline
    \textbf{Sicherheit} & Anfällig für Angriffe & Resistent gegen Padding-Angriffe \\
    \hline
    \textbf{Struktur} & Einfach: $H(k|m)$ & Komplex: zwei Hash-Runden \\
    \hline
    \textbf{Seitenkana.-Schutz} & Nein & Ja (über XOR-Operationen) \\
    \hline
    \textbf{Standardisiert} & Nein & Ja (RFC 2104) \\
    \hline
\end{tabular}

\textbf{(b) Angriffe auf Secret-Prefix und Secret-Suffix:} \checkmark

\textbf{Angriff auf Secret-Suffix $H(m \| k)$:}

\begin{enumerate}[1.]
    \item Längen-Erweiterungs-Angriff (Length Extension Attack)
    \item Blockbasierte Hash-Funktionen (SHA-1, SHA-256) sind anfällig
    \item \textbf{Ablauf:}
    \begin{itemize}
        \item Angreifer kennt: Nachricht $m$, Hash $h = H(m \| k)$ und wahrscheinliche Länge von $k$
        \item Berechnet interner Zustand von SHA nach Verarbeitung von $m \| k$
        \item Kann dann beliebige Daten $m'$ anhängen
        \item Hash für $m \| k \| m'$ berechnen, ohne $k$ zu kennen
    \end{itemize}
    \item \textbf{Beispiel:} Administrator-Kommando $m$ wird mit Secret-Suffix gehashed
    \item Angreifer kann zusätzliche Kommandos anhängen

\end{enumerate}

\end{tcolorbox}
\begin{tcolorbox}[colback=solutionbg, colframe=black!30]

\textbf{Angriff auf Secret-Prefix $H(k \| m)$:}

\begin{enumerate}[1.]
    \item \textbf{Hash-Kollisions-Angriff:}
    \begin{itemize}
        \item Angreifer findet zwei Nachrichten $m_1, m_2$ mit $H(m_1) = H(m_2)$
        \item Aber: $H(k \| m_1) \neq H(k \| m_2)$ im Allgemeinen
    \end{itemize}
    
    \item \textbf{Brute-Force-Angriff:}
    \begin{itemize}
        \item Wenn $k$ kurz ist, kann man alle möglichen $k$ durchprobieren
        \item Überprüfe: Gilt $H(k \| m) = h_{known}$?
        \item Schlüssel-Länge ist der limitierende Faktor
    \end{itemize}
    
    \item \textbf{State-Recovery (bei bestimmten Hashes):}
    \begin{itemize}
        \item Bei manchen Hash-Funktionen kann innerer Zustand rekonstruiert werden
        \item Erlaubt dann Length-Extension-Angriff ähnlich wie Secret-Suffix
    \end{itemize}
\end{enumerate}

\textbf{(c) Warum HMAC-SHA-256 sicherer ist:} \checkmark

\begin{enumerate}[1.]
    \item \textbf{Resistent gegen Längen-Erweiterung:}
    \begin{itemize}
        \item Double-Hashing Struktur verhindert Längen-Erweiterung
        \item Innerer Hash: $H((k \oplus ipad) \| m)$
        \item Äußerer Hash: $H((k \oplus opad) \| innerer\_Hash)$
        \item Angreifer kann nicht einfach $m'$ anhängen
    \end{itemize}
    
    \item \textbf{Schlüssel-Isolation:}
    \begin{itemize}
        \item Schlüssel wird mit ipad und opad verarbeitet (XOR)
        \item Verhindert direkte Verwendung von $k$ im Hash-Prozess
        \item Schlüssel wird vor jedem Hash-Prozess neu
    \end{itemize}
    
    \item \textbf{FIPS-Standard:}
    \begin{itemize}
        \item HMAC ist in RFC 2104 und FIPS 198 standardisiert
        \item Sollte bei allen modernen Protokollen verwendet werden
        \item TLS 1.3, OAuth 2.0, JWT verwenden HMAC
    \end{itemize}
    
    \item \textbf{Praktische Sicherheit:}
    \begin{itemize}
        \item Auch mit nur 256-Bit Schlüssel ist Brute-Force nicht möglich
        \item $2^{256}$ ist astronomisch groß
    \end{itemize}
\end{enumerate}

\end{tcolorbox}

\textbf{Punkteverteilung:}
\begin{itemize}
    \item HMAC vs. einfache Verfahren (3 Punkte) \checkmark
    \item Angriffe auf beide Varianten (3 Punkte) \checkmark
    \item Vorteile HMAC-SHA-256 (3 Punkte) \checkmark
\end{itemize}

% ============================================================================

\newpage

\section{Asymmetrische Verfahren und RSA [17 Punkte]}

\subsection{Aufgabe 4.1 - RSA Grundlagen (10 Punkte)}

\subsubsection*{Musterlösung}

\begin{tcolorbox}[colback=solutionbg, colframe=black!30]

Gegeben: $(e=5, n=21)$, $(d=17)$

\textbf{(a) Verschlüsseln der Nachricht ``5'':} \checkmark

RSA Verschlüsselung: $c = m^e \bmod n$

\begin{align*}
c &= 5^5 \bmod 21 \\
  &= 3125 \bmod 21 \\
  &= 3125 - 148 \times 21 \\
  &= 3125 - 3108 \\
  &= 17
\end{align*}

\textbf{Geheimtext: } $\boxed{17}$

\textbf{(b) Entschlüsseln der Nachricht ``11'':} \checkmark

RSA Entschlüsselung: $m = c^d \bmod n$

\begin{align*}
m &= 11^{17} \bmod 21
\end{align*}

Dies ist schwierig zu rechnen. Wir verwenden Square-and-Multiply:

$17 = 16 + 1 = (10001)_2$

\begin{align*}
11^1 &\equiv 11 \pmod{21} \\
11^2 &\equiv 121 \equiv 16 \pmod{21} \\
11^4 &\equiv 16^2 = 256 \equiv 4 \pmod{21} \\
11^8 &\equiv 4^2 = 16 \pmod{21} \\
11^{16} &\equiv 16^2 = 256 \equiv 4 \pmod{21}
\end{align*}

Daher:
\begin{align*}
11^{17} &= 11^{16} \times 11^1 \\
       &\equiv 4 \times 11 \pmod{21} \\
       &\equiv 44 \pmod{21} \\
       &\equiv 2 \pmod{21}
\end{align*}

\textbf{Klartext: } $\boxed{2}$

Verifikation: $2^5 = 32 = 21 + 11 \equiv 11 \pmod{21}$ \checkmark

\end{tcolorbox}
\begin{tcolorbox}[colback=solutionbg, colframe=black!30]

\textbf{(c) Warum ist dieses Beispiel nicht sicher:} \checkmark

\begin{enumerate}[1.]
    \item \textbf{$n$ ist viel zu klein (21):}
    \begin{itemize}
        \item Schlüsselraum ist nur $\phi(21) = \phi(3 \times 7) = 2 \times 6 = 12$
        \item Brute-Force über alle möglichen Klartexte $< 21$ ist trivial
        \item In der Praxis sollte $n \approx 2048$ Bit sein
    \end{itemize}
    
    \item \textbf{$n$ ist leicht zu faktorisieren:}
    \begin{itemize}
        \item $21 = 3 \times 7$ ist sofort ersichtlich
        \item Wenn ein Angreifer $n$ faktorisieren kann, kann er $\phi(n)$ berechnen
        \item Aus $\phi(n)$ und $e$ kann man $d$ berechnen
    \end{itemize}
    
    \item \textbf{Keine Padding:}
    \begin{itemize}
        \item Praktisches RSA verwendet OAEP-Padding
        \item Verhindert Angriffe wie Textbook-RSA Padding-Oracle
    \end{itemize}
\end{enumerate}

\textbf{(d) Mindestgröße von $n$:} \checkmark

\begin{enumerate}[1.]
    \item \textbf{Aktueller Standard (2024):}
    \begin{itemize}
        \item \textbf{Minimum:} $n \geq 2048$ Bit (empfohlen)
        \item \textbf{Sicher bis ca. 2030}
        \item Entspricht $2^{2048}$ mögliche Werte für $n$
    \end{itemize}
    
    \item \textbf{Zukünftssicherheit:}
    \begin{itemize}
        \item \textbf{Empfohlen:} $n = 4096$ Bit für Langzeit-Sicherheit
        \item \textbf{Post-Quantum:} RSA wird durch Quantencomputer bedroht
        \item Shor's Algorithmus kann RSA in polynomieller Zeit brechen
    \end{itemize}
    
    \item \textbf{Faktorisierungs-Schwierigkeit:}
    \begin{itemize}
        \item Mit aktueller Rechenleistung (2024)
        \item 512 Bit $n$: In wenigen Stunden faktorisierbar
        \item 1024 Bit $n$: Experimentell möglich (mit sehr großem Aufwand)
        \item 2048 Bit $n$: Praktisch unmöglich
    \end{itemize}
\end{enumerate}

\end{tcolorbox}

\textbf{Punkteverteilung:}
\begin{itemize}
    \item Verschlüsselung (2 Punkte) \checkmark
    \item Entschlüsselung (3 Punkte) \checkmark
    \item Sicherheitsprobleme (3 Punkte) \checkmark
    \item Mindestgröße $n$ (2 Punkte) \checkmark
\end{itemize}

\subsection{Aufgabe 4.2 - Digitale Signatur (7 Punkte)}

\subsubsection*{Musterlösung}

\begin{tcolorbox}[colback=solutionbg, colframe=black!30]

\textbf{(a) Funktionsweise einer digitalen Signatur:} \checkmark

\textbf{Generelles Prinzip:}

\begin{enumerate}[\textbf{Schritt 1:}]
    \item \textbf{Hash der Nachricht berechnen:}
    $$h = H(m)$$
    Wobei $H$ eine sichere Hashfunktion ist (z.B. SHA-256)
    
    \item \textbf{Hash mit privatem Schlüssel signieren:}
    $$sig = h^d \bmod n \quad \text{(bei RSA)}$$
    Wobei $d$ der private Schlüssel des Absenders ist
    
    \item \textbf{Nachricht und Signatur versenden:}
    Sender versendet: $(m, sig)$
    
    \item \textbf{Empfänger verifiziert:}
    \begin{itemize}
        \item Hash der erhaltenen Nachricht berechnen: $h' = H(m)$
        \item Signatur mit öffentlichem Schlüssel entschlüsseln: $h_{ver} = sig^e \bmod n$
        \item Vergleichen: Ist $h' = h_{ver}$?
        \item Wenn ja $\to$ Signatur gültig
        \item Wenn nein $\to$ Signatur ungültig oder Nachricht verfälscht
    \end{itemize}
\end{enumerate}

\textbf{Beispiel mit unseren RSA-Parametern:}

Absender möchte Nachricht $m = 5$ signieren. Mit $(d=17, n=21)$:

\begin{align*}
sig &= H(5)^{17} \bmod 21 \\
    &= 5^{17} \bmod 21
\end{align*}

(Annahme: $H(5) = 5$ für dieses Beispiel)

Berechnung (mit $e=5$):
$$sig^e \bmod 21 = sig^5 \bmod 21 = 5$$

Empfänger erhält $(m=5, sig)$ und verifiziert mit öffentlichem Schlüssel $(e=5, n=21)$

\end{tcolorbox}
\begin{tcolorbox}[colback=solutionbg, colframe=black!30]

\textbf{(b) Schutzziele einer digitalen Signatur:} \checkmark

\begin{enumerate}[1.]
    \item \textbf{Authentizität:}
    \begin{itemize}
        \item Nur der Absender (Besitzer des privaten Schlüssels) kann Signatur erzeugen
        \item Empfänger weiß, dass Nachricht wirklich vom Absender kommt
    \end{itemize}
    
    \item \textbf{Integrität:}
    \begin{itemize}
        \item Wenn Nachricht auch nur um ein Bit verändert wird
        \item Ändert sich der Hash komplett (Avalanche-Effekt)
        \item Signatur wird ungültig
        \item Empfänger erkennt sofort Verfälschung
    \end{itemize}
    
    \item \textbf{Verbindlichkeit/Non-Repudiation:}
    \begin{itemize}
        \item Absender kann nicht behaupten, Nachricht nicht geschrieben zu haben
        \item Nur er besitzt privaten Schlüssel
        \item Signatur beweist: Nachricht kam von ihm
        \item Rechtsgültig in vielen Ländern (digitale Signaturen)
    \end{itemize}
    
    \item \textbf{NICHT geschützt: Vertraulichkeit}
    \begin{itemize}
        \item Signatur schützt nicht vor Abhören
        \item Nachricht selbst ist im Klartext (außer zusätzliche Verschlüsselung)
        \item Für Vertraulichkeit: zusätzliche Verschlüsselung notwendig
    \end{itemize}
\end{enumerate}

\end{tcolorbox}
\begin{tcolorbox}[colback=solutionbg, colframe=black!30]

\textbf{(c) Kann Angreifer gültige Signatur erzeugen?} \checkmark

\textbf{Antwort: NEIN (wenn privater Schlüssel geschützt ist)}

\begin{enumerate}[1.]
    \item \textbf{Warum nicht:}
    \begin{itemize}
        \item Signatur erfordert privaten Schlüssel $d$
        \item Nur Absender kennt $d$
        \item Angreifer kennt nur öffentlichen Schlüssel $e$
        \item Aus $e$ kann man nicht auf $d$ schließen (RSA-Annahme)
    \end{itemize}
    
    \item \textbf{Sicherheit basiert auf:}
    \begin{itemize}
        \item \textbf{Faktorisierungs-Annahme:} Schwer, $n = p \times q$ zu faktorisieren
        \item Wenn man $p$ und $q$ kennt $\to$ kann man $\phi(n)$ berechnen
        \item Aus $\phi(n)$ und $e$ $\to$ kann man $d$ berechnen
        \item Aber: Große $n$ sind praktisch nicht faktorisierbar
    \end{itemize}
    
    \item \textbf{Mögliche Angriffe (aber nicht relevant hier):}
    \begin{itemize}
        \item \textbf{Padding-Oracle Attacken:} Bei schwachem Padding möglich
        \item \textbf{Seitenkana-Angriffe:} Durch Timing oder Power-Analysis
        \item \textbf{Quantum-Angriff:} Shor's Algorithmus könnte $n$ faktorisieren
        \item \textbf{Schlüssel-Diebstahl:} Wenn privater Schlüssel gestohlen wird
    \end{itemize}
\end{enumerate}

\end{tcolorbox}

\textbf{Punkteverteilung:}
\begin{itemize}
    \item Funktionsweise erklärt (3 Punkte) \checkmark
    \item Schutzziele (2 Punkte) \checkmark
    \item Kann Angreifer fälschen? (2 Punkte) \checkmark
\end{itemize}

% ============================================================================

\newpage

\section{Authentifikation und Identifikation [15 Punkte]}

\subsection{Aufgabe 5.1 - Passwort-Authentifikation (8 Punkte)}

\subsubsection*{Musterlösung}

\begin{tcolorbox}[colback=solutionbg, colframe=black!30]

\textbf{(a) Identifikation vs. Authentifikation:} \checkmark

\centering
\begin{tabular}{|l|l|l|}
    \hline
    \textbf{Aspekt} & \textbf{Identifikation} & \textbf{Authentifikation} \\
    \hline
    \textbf{Definition} & Wer bist du? & Beweise, dass du wirklich der bist \\
    \hline
    \textbf{In Login:} & Benutzername & Passwort \\
    \hline
    \textbf{Ziel} & Benutzer erkennen & Identität überprüfen \\
    \hline
    \textbf{Beispiel} & ``Mein Name ist Alice'' & ``Mein Passwort ist XYZ123'' \\
    \hline
\end{tabular}

\textbf{Details:}

\begin{itemize}
    \item \textbf{Identifikation:}
    \begin{itemize}
        \item Behauptung einer Identität
        \item Kein Beweis notwendig
        \item Einfach, kann von jedem gemacht werden
        \item Beispiel: ``Ich bin Moritz''
    \end{itemize}
    
    \item \textbf{Authentifikation:}
    \begin{itemize}
        \item Verifizierung, dass die behauptete Identität wahr ist
        \item Beweis der Identität erforderlich
        \item Schwer oder unmöglich zu fälschen
        \item Beispiel: ``Mein Passwort ist [geheim]''
    \end{itemize}
\end{itemize}

\textbf{(b) PAP vs. CHAP:} \checkmark

\centering
\begin{tabular}{|l|l|l|}
    \hline
    \textbf{Kriterium} & \textbf{PAP} & \textbf{CHAP} \\
    \hline
    \textbf{Name} & Password Auth. Prot. & Challenge Handshake Auth. Prot. \\
    \hline
    \textbf{Passwort-Übertragung} & IM KLARTEXT! & Nicht übertragen \\
    \hline
    \textbf{Challenge-Response} & Nein & Ja \\
    \hline
    \textbf{Wiederholung möglich} & Ja (Replay-Angriff!) & Nein (neue Challenge) \\
    \hline
    \textbf{Sicherheit} & SCHWACH & BESSER \\
    \hline
\end{tabular}

\end{tcolorbox}
\begin{tcolorbox}[colback=solutionbg, colframe=black!30]

\textbf{PAP - Password Authentication Protocol:}

\begin{enumerate}[1.]
    \item Client sendet: Benutzername + Passwort (unverschlüsselt!)
    \item Server prüft Passwort
    \item Problem: Passwort kann auf Leitung abgehört werden
\end{enumerate}

\textbf{CHAP - Challenge Handshake Authentication Protocol:}

\begin{enumerate}[1.]
    \item Server sendet: Challenge (Zufallszahl) 
    \item Client berechnet: $response = H(Passwort \| Challenge)$
    \item Client sendet: response (nicht das Passwort!)
    \item Server prüft: Ist $response = H(Passwort \| Challenge)$?
    \item Passwort wird nie übertragen!
\end{enumerate}

\textbf{Sicherheitsimplikationen:} \checkmark

\begin{itemize}
    \item \textbf{PAP Schwächen:}
    \begin{itemize}
        \item Sniffing möglich: Angreifer kann Passwort abhören
        \item Replay möglich: Alte Paket-Sequenzen können erneut gesendet werden
        \item Sollte NIE über unverschlüsselte Verbindung verwendet werden
    \end{itemize}
    
    \item \textbf{CHAP Vorteile:}
    \begin{itemize}
        \item Passwort wird nicht übertragen $\to$ Sniffing funktioniert nicht
        \item Neue Challenge bei jedem Versuch $\to$ Replay unmöglich
        \item Aber: Brute-Force auf Challenge/Response möglich (mit bekanntem Passwort)
    \end{itemize}
    
    \item \textbf{Noch besser:} CHAP über SSL/TLS + RADIUS für zentrale Verwaltung
\end{itemize}

\end{tcolorbox}
\begin{tcolorbox}[colback=solutionbg, colframe=black!30]

\textbf{(c) Passwort-Hash mit Salt vs. Klartext:} \checkmark

\textbf{Warum Hash+Salt sicherer ist:}

\begin{enumerate}[1.]
    \item \textbf{Klartext-Passwörter sind kritisch:}
    \begin{itemize}
        \item Wenn Datenbank gestohlen wird $\to$ alle Passwörter sind kompromittiert
        \item Sofort verwendet werden können (z.B. bei anderen Services)
        \item Keine Chance, Benutzer zu warnen
    \end{itemize}
    
    \item \textbf{Hash-Vorteile:}
    \begin{itemize}
        \item Passwort wird mit Hashfunktion verschlüsselt (One-Way)
        \item Wenn Datenbank gestohlen wird: nur Hashes sind sichtbar
        \item Angreifer müsste Passwörter eraten und hashen
        \item Brute-Force-Angriff wird teuer
    \end{itemize}
    
    \item \textbf{Salt-Vorteile:}
    \begin{itemize}
        \item Zufallswert, der mit Passwort konkateniert wird
        \item Vor dem Hashen: $H(salt \| password)$
        \item Verhindert Rainbow-Tables: vorgefertigte Hash-Tabellen
        \item Auch identische Passwörter haben unterschiedliche Hashes
    \end{itemize}
    
    \item \textbf{Beispiel - ohne Salt:}
    \begin{itemize}
        \item Benutzer 1 Passwort: ``123456'' $\to$ Hash: abc123
        \item Benutzer 2 Passwort: ``123456'' $\to$ Hash: abc123
        \item Angreifer sieht: Viele Benutzer haben gleichen Hash
        \item Mit Rainbow-Table: Quick-Match gegen bekannte Passwörter
    \end{itemize}
    
    \item \textbf{Beispiel - mit Salt:}
    \begin{itemize}
        \item Benutzer 1: salt=xyz, password=``123456'' $\to$ $H(\text{xyz123456})$ = abc123
        \item Benutzer 2: salt=abc, password=``123456'' $\to$ $H(\text{abc123456})$ = def456
        \item Angreifer sieht: Unterschiedliche Hashes
        \item Rainbow-Table ist nutzlos (müsste alle Salts kennen)
    \end{itemize}
    
    \item \textbf{Moderne Best-Practice:}
    \begin{itemize}
        \item Verwende: bcrypt, scrypt oder Argon2
        \item Diese haben automatisch Salt und sind langsam (Cost-Faktor)
        \item Brute-Force wird dadurch extrem teuer
    \end{itemize}
\end{enumerate}

\end{tcolorbox}

\textbf{Punkteverteilung:}
\begin{itemize}
    \item Identifikation vs. Authentifikation (2 Punkte) \checkmark
    \item PAP vs. CHAP Unterschiede (3 Punkte) \checkmark
    \item Hash+Salt Sicherheit (3 Punkte) \checkmark
\end{itemize}

\subsection{Aufgabe 5.2 - Mehfaktor-Authentifikation (7 Punkte)}

\subsubsection*{Musterlösung}

\begin{tcolorbox}[colback=solutionbg, colframe=black!30]

\textbf{(a) Die drei Authentifikationsfaktoren:} \checkmark

\begin{enumerate}[1.]
    \item \textbf{Wissen (Something you know):}
    \begin{itemize}
        \item Passwort, PIN, Sicherheitsfrage
        \item Nur die Person sollte es wissen
    \end{itemize}
    
    \item \textbf{Besitz (Something you have):}
    \begin{itemize}
        \item Smartphone, Hardware-Token, SmartCard
        \item Physisches Gerät wird benötigt
    \end{itemize}
    
    \item \textbf{Biometrie (Something you are):}
    \begin{itemize}
        \item Fingerabdruck, Gesichtserkenn, Iris-Scan
        \item Körperliche/biologische Eigenschaft
    \end{itemize}
\end{enumerate}

\textbf{Zusätzliche Faktoren (Erwähnung möglich):}
\begin{itemize}
    \item \textbf{Standort (Where you are):} Geografischer Ort
    \item \textbf{Zeit (When you are):} Zeitlich-basierte OTP
\end{itemize}

\textbf{(b) Beispiele für Zweifaktor-Authentifikation (2FA):} \checkmark

\textbf{Beispiel 1: Passwort + SMS-Code}
\begin{enumerate}[1.]
    \item Faktoren: Wissen (Passwort) + Besitz (Smartphone)
    \item Ablauf:
    \begin{itemize}
        \item Benutzer gibt Passwort ein
        \item System versendet SMS mit 6-stelligem Code
        \item Benutzer muss Code eingeben
        \item Benutzer ist authentifiziert
    \end{itemize}
    \item Vorteile:
    \begin{itemize}
        \item Weit verbreitet und einfach
        \item Funktioniert auf jedem Telefon
        \item Günstig umzusetzen
    \end{itemize}
    \item Nachteile:
    \begin{itemize}
        \item SMS kann abgefangen werden (SIM-Swapping)
        \item Verzögerung beim SMS-Versand
    \end{itemize}
\end{enumerate}

\end{tcolorbox}
\begin{tcolorbox}[colback=solutionbg, colframe=black!30]

\textbf{Beispiel 2: Passwort + Authenticator-App (TOTP)}
\begin{enumerate}[1.]
    \item Faktoren: Wissen (Passwort) + Biometrie/Besitz (Fingerprint auf Handy)
    \item Ablauf:
    \begin{itemize}
        \item App wie Google Authenticator oder Authy
        \item Generiert zeitbasierte 6-stellige Codes (alle 30 Sek.)
        \item QR-Code wird beim Setup gescannt (Shared Secret)
        \item Benutzer gibt Code ein
    \end{itemize}
    \item Vorteile:
    \begin{itemize}
        \item Funktioniert offline (kein Internetzugang nötig)
        \item Sicherer als SMS (kein physischer Übertragungskanal)
        \item Keine monatlichen SMS-Kosten
        \item Backup-Codes möglich
    \end{itemize}
    \item Nachteile:
    \begin{itemize}
        \item Benutzer muss App installieren und verwenden
        \item Bei Telefon-Verlust: Zugriff verloren (ohne Backup-Codes)
    \end{itemize}
\end{enumerate}

\textbf{Weitere Beispiele:}
\begin{itemize}
    \item \textbf{Passwort + Hardware-Token:} z.B. YubiKey
    \item \textbf{Passwort + Fingerprint:} z.B. Windows Hello, iPhone FaceID
    \item \textbf{Passwort + Email-Bestätigung:} Weniger sicher, aber immer noch 2FA
\end{itemize}

\textbf{(c) Welche Angriffe werden durch MFA NICHT verhindert?} \checkmark

\begin{enumerate}[1.]
    \item \textbf{Phishing-Angriffe (teilweise):}
    \begin{itemize}
        \item Benutzer wird auf gefälschte Login-Seite gelockt
        \item Gibt Passwort + SMS-Code ein (einfaches MFA)
        \item Hacker kann damit sofort einloggen
        \item Besser: Phishing-resistant MFA (FIDO2/WebAuthn)
    \end{itemize}
    
    \item \textbf{Malware auf dem Endgerät:}
    \begin{itemize}
        \item Keystroke-Logging: alle Eingaben werden aufgezeichnet
        \item Screen-Scraping: Screenshots zeigen SMS-Codes
        \item MFA hilft nicht, wenn beide Faktoren auf gleich-Gerät
    \end{itemize}
\end{enumerate}

\end{tcolorbox}
\begin{tcolorbox}[colback=solutionbg, colframe=black!30]
    
\begin{enumerate}[3.]
    \item \textbf{Man-in-the-Middle (bei schwachem MFA):}
    \begin{itemize}
        \item Angreifer positioniert sich zwischen Benutzer und Server
        \item Fängt Passwort + MFA-Code ab
        \item Sendet diese sofort weiter
        \item Benutzer wundert sich über mehrere Login-Versuche
    \end{itemize}
    
    \item \textbf{SIM-Swapping (bei SMS-MFA):}
    \begin{itemize}
        \item Angreifer überzeugt Telefonprovider, SIM zu wechseln
        \item SMS-Codes gehen jetzt an Angreifer
        \item MFA ist nutzlos
    \end{itemize}
    
    \item \textbf{Session-Hijacking:}
    \begin{itemize}
        \item Nach erfolgreicher Authentifikation: Session-Cookie wird gestohlen
        \item Angreifer kann mit Cookie weitermachen
        \item MFA war nur beim Login relevant
        \item Lösung: Kontinuierliche Verifizierung
    \end{itemize}
    
    \item \textbf{Sozial-Engineering:}
    \begin{itemize}
        \item Angreifer redet Benutzer um, Zugang zu geben
        \item z.B. ``Ich bin vom Support, bitte gib mir deinen Code''
        \item MFA hilft nicht gegen menschliche Manipulation
    \end{itemize}
    
    \item \textbf{Kompromittierter Server:}
    \begin{itemize}
        \item Wenn der Server gehackt wurde
        \item MFA schützt den Server selbst nicht
        \item Daten auf dem Server können gestohlen werden
    \end{itemize}
\end{enumerate}

\textbf{Beste MFA-Variante gegen Phishing:}

\begin{itemize}
    \item \textbf{FIDO2/WebAuthn:} Phishing-resistent
    \begin{itemize}
        \item Hardware-Key oder Biometrie
        \item Verifiziert die Domain
        \item Phishing-Seite wird nicht funktionieren (falsche Domain)
    \end{itemize}
\end{itemize}

\end{tcolorbox}

\textbf{Punkteverteilung:}
\begin{itemize}
    \item Drei Faktoren genannt (2 Punkte) \checkmark
    \item Zwei 2FA-Beispiele mit Vorteilen (3 Punkte) \checkmark
    \item Angriffe, die MFA nicht verhindert (2 Punkte) \checkmark
\end{itemize}

% ============================================================================

\newpage

\section{DSGVO und Security Engineering [15 Punkte]}

\subsection{Aufgabe 6.1 - DSGVO (10 Punkte)}

\subsubsection*{Musterlösung}

\begin{tcolorbox}[colback=solutionbg, colframe=black!30]

\textbf{(a) Definition ``personenbezogene Daten'':} \checkmark

Nach Artikel 4 DSGVO:

\textit{``Personenbezogene Daten sind alle Informationen, die sich auf eine identifizierte oder identifizierbare natürliche Person (betroffene Person) beziehen.''}

\textbf{Details:}

\begin{itemize}
    \item Direkte Identifikatoren: Name, Adresse, Ausweisnummer
    \item Indirekte Identifikatoren: IP-Adresse, Cookie-ID, Kundennummer
    \item Pseudonymisierte Daten: Noch personenbezogen, aber nicht direkt identifizierbar
    \item Anonyme Daten: Nicht personenbezogen (wenn wirklich nicht zurückverfolgbar)
    \item Online-Identifikatoren: Login-Namen, E-Mail-Adressen
\end{itemize}

\textbf{Wichtig:} Auch ``sensible'' Daten (z.B. Telefonnummer, Standort) können personenbezogen sein, auch wenn sie nicht ``special category'' sind.

\textbf{(b) Daten besonderer Kategorie (Artikel 9):} \checkmark

Mindestens 4 Kategorien:

\begin{enumerate}[1.]
    \item \textbf{Rassische oder ethnische Herkunft}
    \begin{itemize}
        \item Beispiele: Ethnisches Profiling, Diskriminierung
    \end{itemize}
    
    \item \textbf{Politische Meinungen}
    \begin{itemize}
        \item Beispiele: Wahl-Verhalten, Parteimitgliedschaft
    \end{itemize}
    
    \item \textbf{Religiöse oder weltanschauliche Überzeugungen}
    \begin{itemize}
        \item Beispiele: Glaubensgemeinschaft, Kirche
    \end{itemize}
    
    \item \textbf{Gewerkschaftszugehörigkeit}
    \begin{itemize}
        \item Beispiele: Mitgliedschaft in Gewerkschaft
    \end{itemize}
    
    \item \textbf{Genetische Daten}
    \begin{itemize}
        \item Beispiele: DNA-Profil, genetische Testergebnisse
    \end{itemize}
    
    \item \textbf{Biometrische Daten}
    \begin{itemize}
        \item Beispiele: Fingerabdruck, Gesichtserkennungs-Daten, Iris-Scan
        \item Zur Identifikation einer Person
    \end{itemize}
    
    \item \textbf{Gesundheitsdaten}
    \begin{itemize}
        \item Beispiele: Medizinische Diagnose, Behandlung, Rezepte
    \end{itemize}
    
    \item \textbf{Daten zum Sexualleben oder der sexuellen Orientierung}
    \begin{itemize}
        \item Beispiele: Sexuelle Orientierung, Sexualverhalten
    \end{itemize}
\end{enumerate}

\end{tcolorbox}
\begin{tcolorbox}[colback=solutionbg, colframe=black!30]

\textbf{(c) Benutzer speichert Gesichtserkennungs-Daten - welche DSGVO-Artikel:} \checkmark

\begin{enumerate}[1.]
    \item \textbf{Artikel 9 (Besondere Kategorien):}
    \begin{itemize}
        \item Biometrische Daten zur Identifikation fallen unter Art. 9
        \item Verarbeitung ist grundsätzlich verboten!
        \item Nur mit ausdrücklicher Zustimmung oder Ausnahmeregelung
    \end{itemize}
    
    \item \textbf{Artikel 4 (Definition personenbezogene Daten):}
    \begin{itemize}
        \item Gesichtserkennungs-Daten sind personenbezogen
        \item Können einer Person zugeordnet werden
    \end{itemize}
    
    \item \textbf{Artikel 5 (Datenschutzprinzipien):}
    \begin{itemize}
        \item Rechtmäßigkeit, Fairness, Transparenz
        \item Zweckbindung (warum werden Biometrien gespeichert?)
        \item Speicherbegrenzung (wie lange?)
        \item Datenminimierung (sind wirklich Biometrien nötig?)
    \end{itemize}
    
    \item \textbf{Artikel 6 (Rechtmäßigkeit der Verarbeitung):}
    \begin{itemize}
        \item Muss eine rechtliche Grundlage haben
        \item z.B. Zustimmung (Artikel 7), Vertrag, rechtliche Pflicht
    \end{itemize}
    
    \item \textbf{Artikel 21 (Widerspruchsrecht):}
    \begin{itemize}
        \item Betroffene können Verarbeitung ihrer Biometrien widersprechen
    \end{itemize}
    
    \item \textbf{Artikel 32 (Datenschutz durch Technik und Organisati-Org):}
    \begin{itemize}
        \item Besondere Schutzmaßnahmen für biometrische Daten
        \item Verschlüsselung, Zugriffskontrolle, Auditing
    \end{itemize}
\end{enumerate}

\end{tcolorbox}
\begin{tcolorbox}[colback=solutionbg, colframe=black!30]

\textbf{(d) Strafen bei DSGVO-Verstößen:} \checkmark

\textbf{Verwaltungsgeldstrafen nach Artikel 83:}

\begin{enumerate}[1.]
    \item \textbf{Kategoriestufe 1 (bis zu 10.000.000 EUR oder 2\% Umsatz):}
    \begin{itemize}
        \item Verletzung von Dokumentationspflichten
        \item Unzureichende Transparenzmaßnahmen
        \item Verletzung technischer Schutzmaßnahmen
    \end{itemize}
    
    \item \textbf{Kategoriestufe 2 (bis zu 20.000.000 EUR oder 4\% Umsatz):}
    \begin{itemize}
        \item Verletzung von Grundprinzipien (Art. 5, 6)
        \item Betroffenenrechte nicht gewährt (Art. 12-22)
        \item Datenschutz-Folgenabschätzung nicht durchgeführt
        \item Sicherheitsverletzung nicht gemeldet
    \end{itemize}
\end{enumerate}

\textbf{Beispiel-Strafen:}

\begin{itemize}
    \item \textbf{Google (2019):} 50 Mio. EUR - Mangelnde Transparenz
    \item \textbf{Amazon (2021):} 746 Mio. EUR - Unrechtmäßige Verarbeitung
    \item \textbf{Meta/Facebook (2022):} 405 Mio. EUR - Datenübermittlung in USA
\end{itemize}

\textbf{Zusätzlich:}

\begin{itemize}
    \item \textbf{Schadensersatz:} Geschädigte können Schadensersatz fordern
    \item \textbf{Strafrechtliche Verfolgung:} In schweren Fällen
    \item \textbf{Reputational Damage:} Vertrauen-Verlust, Kundenverlust
\end{itemize}

\end{tcolorbox}

\textbf{Punkteverteilung:}
\begin{itemize}
    \item Personenbezogene Daten definiert (2 Punkte) \checkmark
    \item 4+ Kategorien benannt (2 Punkte) \checkmark
    \item Relevante DSGVO-Artikel (3 Punkte) \checkmark
    \item Strafen (3 Punkte) \checkmark
\end{itemize}

\subsection{Aufgabe 6.2 - Security Engineering (5 Punkte)}

\subsubsection*{Musterlösung}

\begin{tcolorbox}[colback=solutionbg, colframe=black!30]

\textbf{(a) ``Secure by Design'' Konzept:} \checkmark

\textit{``Secure by Design''} bedeutet, dass Sicherheit von Anfang an in die Entwicklung von Systemen und Produkten integriert wird, nicht erst nachträglich hinzugefügt.

\textbf{Detaillierte Erklärung:}

\begin{enumerate}[1.]
    \item \textbf{Sicherheit als primäres Ziel:}
    \begin{itemize}
        \item Schon bei Anforderungsdefinition Sicherheitsanforderungen festlegen
        \item Nicht: ``Erst machen, dann sichern''
        \item Sondern: ``Von Anfang an sicher''
    \end{itemize}
    
    \item \textbf{Architektur-Design:}
    \begin{itemize}
        \item Sicherheit in der Systemarchitektur berücksichtigen
        \item z.B. Least Privilege, Defense in Depth, Zero Trust
        \item Modulare Architektur für Isolation kritischer Komponenten
    \end{itemize}
    
    \item \textbf{Threat Modeling:}
    \begin{itemize}
        \item Frühzeitig potenzielle Angriffe identifizieren
        \item STRIDE-Methode oder PASTA
        \item Gegenmassnahmen vorausplanen
    \end{itemize}
    
    \item \textbf{Sichere Kodierungsstandards:}
    \begin{itemize}
        \item OWASP Top 10 vermeiden
        \item Code-Reviews und Static Analysis
        \item Sichere Bibliotheken und Frameworks verwenden
    \end{itemize}
    
    \item \textbf{Testing und Validierung:}
    \begin{itemize}
        \item Security-Tests in jedem Development-Sprint
        \item Penetration Testing vor Release
        \item Vulnerability Scanning automatisiert
    \end{itemize}
\end{enumerate}

\end{tcolorbox}
\begin{tcolorbox}[colback=solutionbg, colframe=black!30]

\textbf{Beispiel:}

\begin{itemize}
    \item \textbf{Schlecht (Ad-hoc Sicherheit):}
    \begin{itemize}
        \item Entwickler bauen Web-App: ``Später kümmern wir uns um Sicherheit''
        \item Nach 1 Jahr: SQL-Injection Lücken entdeckt
        \item Nachträgliche Fixes sind teuer und fehleranfällig
    \end{itemize}
    
    \item \textbf{Gut (Secure by Design):}
    \begin{itemize}
        \item Vor Entwicklung: Threat Model erstellt
        \item Parameter-Validierung eingeplant
        \item ORM-Framework gewählt (gegen SQL-Injection)
        \item Code-Review mit Security-Team
        \item Keine Sicherheitslücken von Anfang an
    \end{itemize}
\end{itemize}

\textbf{(b) Zwei Maßnahmen für ``Secure by Default'':} \checkmark

\textbf{Maßnahme 1: Sichere Standard-Konfiguration}

\begin{itemize}
    \item \textbf{Definition:} System startet mit maximaler Sicherheit, nicht minimaler
    \item \textbf{Beispiele:}
    \begin{itemize}
        \item Passwort-Komplexitäts-Anforderungen sind Standard (nicht optional)
        \item TLS/HTTPS ist Standard (nicht HTTP)
        \item Zwei-Faktor-Authentifikation ist enabled by default
        \item Firewall ist aktiviert (nicht disabled)
        \item Datenbank-Verbindungen werden verschlüsselt (nicht im Klartext)
    \end{itemize}
    \item \textbf{Vorteil:} Auch wenn Administrator nicht aufpasst, ist System sicher
    \item \textbf{Benutzer müssen aktiv Sicherheit deaktivieren (wenn sie es wollen)}
    \begin{itemize}
        \item Anstatt: Benutzer müssen aktiv Sicherheit aktivieren
        \item Was die meisten vergessen
    \end{itemize}
\end{itemize}

\end{tcolorbox}
\begin{tcolorbox}[colback=solutionbg, colframe=black!30]

\textbf{Maßnahme 2: Minimales Privilege-Level (Principle of Least Privilege)}

\begin{itemize}
    \item \textbf{Definition:} Benutzer/Prozesse bekommen nur minimal notwendige Rechte
    \item \textbf{Beispiele:}
    \begin{itemize}
        \item Web-Server läuft nicht als root, sondern als ``www-data'' User
        \item Datenbank-Account hat nur SELECT-Rechte, nicht DROP/DELETE
        \item Prozesse in Container isoliert (nicht mit Host-Zugriff)
        \item Browser-Extensions nur minimale Permissions (nicht ``alle Websites''
    \end{itemize}
    \item \textbf{Vorteil:} Wenn eine Komponente gehackt wird, ist Schaden begrenzt
    \item \textbf{Beispiel:}
    \begin{itemize}
        \item Hacker kompromittiert Web-App
        \item App läuft als User ``www-data'' mit nur SELECT-Rechten
        \item Hacker kann Datenbank nicht löschen
        \item System bleibt verfügbar
    \end{itemize}
\end{itemize}

\textbf{Weitere mögliche Maßnahmen:}

\begin{itemize}
    \item Daten-Verschlüsselung at-rest und in-transit (Standard)
    \item Audit-Logging aktiviert by default
    \item Automatische Security-Updates
    \item Rate-Limiting gegen Brute-Force (Standard)
    \item Content Security Policy (CSP) Headers
    \item Input-Validation bei allen Eingaben
\end{itemize}

\end{tcolorbox}

\textbf{Punkteverteilung:}
\begin{itemize}
    \item Secure by Design erklärt (2 Punkte) \checkmark
    \item Sichere Standard-Konfiguration (1,5 Punkte) \checkmark
    \item Least Privilege (1,5 Punkte) \checkmark
\end{itemize}

% ============================================================================

\newpage

\section{Praxisaufgabe: Sicherheitsanalyse [10 Punkte]}

\subsection{Aufgabe 7 - Szenarien-Bewertung}

\subsubsection*{Szenario A: Cloud-Authentifikation}

\begin{tcolorbox}[colback=solutionbg, colframe=black!30]

\textbf{Beschreibung:} SaaS-Anbieter mit AES-256-CBC Verschlüsselung, Cloud-Provider Schlüsselverwaltung, Passwort + SMS 2FA

\textbf{(i) Welche Schutzziele sind abgedeckt?} \checkmark

\begin{enumerate}[1.]
    \item \textbf{Vertraulichkeit:}
    \begin{itemize}
        \item AES-256-CBC bietet gute Verschlüsselung
        \item Daten in der Cloud sind verschlüsselt
        \item \textbf{Abgedeckt} \checkmark
    \end{itemize}
    
    \item \textbf{Authentizität:}
    \begin{itemize}
        \item Passwort + 2FA (SMS) bietet Authentifikation
        \item Benutzer kann sich identifizieren
        \item \textbf{Abgedeckt} \checkmark
    \end{itemize}
    
    \item \textbf{Integrität:}
    \begin{itemize}
        \item AES-256-CBC schützt nur vor Abhören, nicht vor Tampering
        \item GCM würde Integrität schützen, aber nur CBC verwendet
        \item \textbf{Nur teilweise abgedeckt}
    \end{itemize}
    
    \item \textbf{Verfügbarkeit:}
    \begin{itemize}
        \item Cloud-Provider kann DDoS-Schutz bieten
        \item Aber nicht aus der Beschreibung ersichtlich
        \item \textbf{Nicht klar} ?
    \end{itemize}
    
    \item \textbf{Verbindlichkeit:}
    \begin{itemize}
        \item 2FA mit Passwort + SMS ermöglicht Attribution
        \item Benutzer kann nicht behaupten, es war nicht er
        \item \textbf{Abgedeckt} \checkmark
    \end{itemize}
\end{enumerate}

\end{tcolorbox}
\begin{tcolorbox}[colback=solutionbg, colframe=black!30]

\textbf{(ii) Welche Risiken bestehen noch?} \checkmark

\begin{enumerate}[1.]
    \item \textbf{Keine Authentizität für Integrität:}
    \begin{itemize}
        \item AES-CBC ist nicht authentifiziert
        \item Wenn Daten manipuliert werden, wird das nicht erkannt
        \item Lösung: AES-GCM oder CBC mit HMAC verwenden
    \end{itemize}
    
    \item \textbf{SMS-MFA ist nicht Phishing-resistent:}
    \begin{itemize}
        \item Phishing-Angreifer kann Passwort + Code abfangen
        \item Wenn Benutzer auf gefälschte Seite geht und beides eingibt
        \item Und Angreifer sendet weiter, ist 2FA nutzlos
    \end{itemize}
    
    \item \textbf{SMS-MFA kann abgefangen werden:}
    \begin{itemize}
        \item SIM-Swapping möglich
        \item Oder Interception auf Telefonprovider-Seite
        \item Besser: TOTP (Time-based OTP) oder FIDO2
    \end{itemize}
    
    \item \textbf{Schlüsselverwaltung durch Cloud-Provider:}
    \begin{itemize}
        \item ``Shared Responsibility Model'' - sind Schlüssel wirklich sicher?
        \item Cloud-Provider könnte Zugriff auf Schlüssel haben
        \item Insider-Angriff möglich
        \item GDPR-Compliance fraglich
    \end{itemize}
    
    \item \textbf{Daten im Transit:}
    \begin{itemize}
        \item Ist Kommunikation zwischen Client und Cloud verschlüsselt?
        \item TLS sollte verwendet werden (nicht nur Daten-Verschlüsselung)
    \end{itemize}
    
    \item \textbf{Keine Daten-Minimierung:}
    \begin{itemize}
        \item Werden wirklich alle Daten verschlüsselt?
        \item Oder gibt es Audit-Logs, die im Klartext sind?
    \end{itemize}
\end{enumerate}

\end{tcolorbox}
\begin{tcolorbox}[colback=solutionbg, colframe=black!30]

\textbf{(iii) Wie würden Sie die Sicherheit verbessern?} \checkmark

\begin{enumerate}[1.]
    \item \textbf{AES-GCM statt AES-CBC:}
    \begin{itemize}
        \item GCM bietet Authentizität + Verschlüsselung in einem
        \item Integritäts-Checks sind automatisch
    \end{itemize}
    
    \item \textbf{Phishing-resistente MFA:}
    \begin{itemize}
        \item FIDO2/WebAuthn mit Hardware-Keys
        \item Oder: Biometrische Authentifikation (Face/Fingerprint)
        \item Zusätzlich: Passwortlos gehen (Passwordless Authentication)
    \end{itemize}
    
    \item \textbf{Bessere Schlüsselverwaltung:}
    \begin{itemize}
        \item Customer-Managed Keys (CMK) statt Provider-Managed Keys
        \item Selbst für Schlüsselhilfsmittel verantwortlich
        \item Key Rotation implementieren
    \end{itemize}
    
    \item \textbf{Zero-Trust Architecture:}
    \begin{itemize}
        \item Assume breach - alle Verbindungen überprüfen
        \item Micro-segmentation der Cloud-Infrastruktur
        \item Continuous Verification statt einmaliger Login
    \end{itemize}
    
    \item \textbf{Weitere Maßnahmen:}
    \begin{itemize}
        \item TLS 1.3 für Datenverkehr
        \item DDoS-Schutz implementieren
        \item Backup und Disaster Recovery
        \item Regelmäßige Penetration-Tests
        \item Security Audit-Logs
    \end{itemize}
\end{enumerate}

\end{tcolorbox}

\subsubsection*{Szenario B: IoT-Geräte}

\begin{tcolorbox}[colback=solutionbg, colframe=black!30]

\textbf{Beschreibung:} Smart-Türschloss mit pre-shared keys in Firmware

\textbf{(i) Identifizieren Sie die Sicherheitsprobleme:} \checkmark

\begin{enumerate}[1.]
    \item \textbf{KRITISCH: Pre-Shared Keys in Firmware:}
    \begin{itemize}
        \item Firmware kann durch Reverse-Engineering extrahiert werden
        \item Alle Türschlösser mit gleicher Firmware haben gleichen Key
        \item Ein gehacktes Gerät $\to$ Alle Geräte gefährdet
        \item Benutzer kann Key nicht ändern
    \end{itemize}
    
    \item \textbf{KRITISCH: Kein Authentisierung der App:}
    \begin{itemize}
        \item Wie wird sichergestellt, dass es die ``richtige'' App ist?
        \item Reverse-Engineering der App offenbart Schlüssel
        \item Gefälschte App könnte Schloss öffnen
    \end{itemize}
    
    \item \textbf{KRITISCH: Keine gegenseitige Authentifikation:}
    \begin{itemize}
        \item Schloss vertraut App automatisch (mit richtigem Key)
        \item Aber: App vertraut Schloss automatisch?
        \item Man-in-the-Middle möglich
    \end{itemize}
    
    \item \textbf{Schwach: Physische Sicherheit:}
    \begin{itemize}
        \item Firmware könnte durch Hardware-Hacking extrahiert werden
        \item Side-Channel Angriffe möglich (Power-Analysis)
        \item Schloss könnte direkt gehackt werden
    \end{itemize}
    
    \item \textbf{Schwach: Keine Replay-Protection:}
    \begin{itemize}
        \item Wenn gleicher Key verwendet wird
        \item Angreifer könnte alte Nachrichten erneut senden
        \item Locking Commands könnten Protolyse sein
    \end{itemize}
    
    \item \textbf{Schwach: Keine Updates möglich:}
    \begin{itemize}
        \item Wenn Sicherheitslücke entdeckt wird
        \item Firmware kann nicht aktualisiert werden
        \item Schlüssel ist für immer in Hardware graviert
    \end{itemize}
\end{enumerate}

\end{tcolorbox}
\begin{tcolorbox}[colback=solutionbg, colframe=black!30]

\textbf{(ii) Wie würde ein Angreifer vorgehen?} \checkmark

\begin{enumerate}[\textbf{Szenario 1: Firmware-Extraktion}]
    \item Angreifer kauft einen Smart-Lock
    \item Öffnet das Gerät physisch
    \item Verbindet sich mit JTAG/Serial-Port
    \item Dumpet die Firmware
    \item Reverse-Engineert die Firmware
    \item Findet den Pre-Shared Key
    \item Kann jetzt alle Schlösser mit dieser Firmware öffnen
\end{enumerate}

\begin{enumerate}[\textbf{Szenario 2: Reverse-Engineering der App}]
    \item Angreifer installiert die legale Smart-Lock App
    \item Dekompiliert die APK (Android) oder Extracts IPA (iOS)
    \item Sucht nach hardcodierten Schlüsseln
    \item Findet Pre-Shared Key in den App-Ressourcen
    \item Schreibt eigene App mit dem Key
    \item Kann jede Tür öffnen
\end{enumerate}

\begin{enumerate}[\textbf{Szenario 3: Funk-Sniffing}]
    \item Angreifer mit SDR (Software Defined Radio)
    \item Benutzer öffnet Schloss über App
    \item Angreifer fängt Funk-Signal ab
    \item Speichert das Unlock-Kommando
    \item Sendet es später erneut (Replay-Angriff)
    \item Schloss öffnet sich
\end{enumerate}

\end{tcolorbox}
\begin{tcolorbox}[colback=solutionbg, colframe=black!30]

\textbf{(iii) Welche Lösung würden Sie empfehlen?} \checkmark

\begin{enumerate}[1.]
    \item \textbf{Asymmetrische Kryptographie (PKI):}
    \begin{itemize}
        \item Jeden Smart-Lock ein einzigartiges Zertifikat geben
        \item Ermöglicht Benutzer-spezifische Schlüssel
        \item Public-Key ist im Schloss, privater Key bei Benutzer
        \item Skalierbar auf neue Benutzer
    \end{itemize}
    
    \item \textbf{Challenge-Response Authentifikation:}
    \begin{itemize}
        \item Nicht einfach Secret Key senden
        \item Schloss sendet Challenge
        \item App antwortet mit: $response = H(Challenge | Secret)$
        \item Verhindert Replay-Angriffe
        \item Challenge sollte Timestamp enthalten
    \end{itemize}
    
    \item \textbf{OTA Updates (Over-The-Air):}
    \begin{itemize}
        \item Firmware regelmäßig aktualisieren
        \item Neue Sicherheits-Patches einspielen
        \item Verschlüsselte Updates übertragen
    \end{itemize}
    
    \item \textbf{Separate Credentials pro Benutzer:}
    \begin{itemize}
        \item Nicht ein Schlüssel für alle
        \item Benutzer bekommt eindeutige Credentials
        \item Ermöglicht Zugriffskontrolle und Revocation
        \item Wenn Benutzer geht: nur seinen Key revoken, nicht für alle
    \end{itemize}
    
    \item \textbf{Sichere Schlüsselspeicherung:}
    \begin{itemize}
        \item TPM (Trusted Platform Module) im Schloss
        \item Verhindert Firmware-Dump
        \item Hardware-basierte Sicherheit
    \end{itemize}
    
    \item \textbf{Verschlüsselte Kommunikation:}
    \begin{itemize}
        \item TLS/DTLS für Funk-Kommunikation
        \item Nicht im Klartext senden
        \item Authentifizierung beider Seiten
    \end{itemize}
\end{enumerate}
    
\end{tcolorbox}
\begin{tcolorbox}[colback=solutionbg, colframe=black!30]
    
\begin{enumerate}[7.]
    \item \textbf{Beispiel-Implementierung:}
    \begin{itemize}
        \item Benutzer registriert sich in App/Web-Portal
        \item Erhält eindeutige Device-ID und Private Key
        \item Smart-Lock hat Public Key des Herstellers und Device-Zertifikat
        \item Beim Öffnen:
        \begin{enumerate}[(1)]
            \item App sendet User-ID + Timestamp
            \item Schloss sendet Challenge + nonce
            \item App signiert: $sig = Sign(Challenge | nonce | Timestamp, PrivateKey)$
            \item Schloss verifiziert mit Public Key
            \item Entsperrt nur wenn alles korrekt
        \end{enumerate}
    \end{itemize}
\end{enumerate}

\end{tcolorbox}

\textbf{Punkteverteilung für Aufgabe 7:}
\begin{itemize}
    \item Szenario A - Schutzziele (2 Punkte) \checkmark
    \item Szenario A - Risiken (2 Punkte) \checkmark
    \item Szenario A - Verbesserungen (1 Punkt) \checkmark
    \item Szenario B - Sicherheitsprobleme (2 Punkte) \checkmark
    \item Szenario B - Angreifer-Vorgehen (1 Punkt) \checkmark
    \item Szenario B - Lösungsvorschlag (2 Punkte) \checkmark
\end{itemize}

% ============================================================================

\newpage

\section*{ABSCHLIESSENDE HINWEISE}

\subsection*{Gesamtpunkte und Gewichtung}

Die maximale Punktzahl ist 110 Punkte. Die Verteilung nach Aufgaben:

\centering
\begin{tabular}{|c|c|c|}
    \hline
    \textbf{Aufgabe} & \textbf{Punkte} & \textbf{Prozent} \\
    \hline
    1 & 15 & 13,6\% \\
    2 & 20 & 18,2\% \\
    3 & 18 & 16,4\% \\
    4 & 17 & 15,5\% \\
    5 & 15 & 13,6\% \\
    6 & 15 & 13,6\% \\
    7 & 10 & 9,1\% \\
    \hline
    \textbf{TOTAL} & \textbf{110} & \textbf{100\%} \\
    \hline
\end{tabular}

\subsection*{Bewertungsskala}

\centering
\begin{tabular}{|c|c|c|c|}
    \hline
    \textbf{Note} & \textbf{Punkte} & \textbf{Prozent} & \textbf{Beschreibung} \\
    \hline
    1,0 & $\geq 99$ & $\geq 90\%$ & Ausgezeichnet \\
    1,3 & $\geq 90$ & $\geq 82\%$ & Sehr gut \\
    1,7 & $\geq 82$ & $\geq 75\%$ & Gut \\
    2,0 & $\geq 77$ & $\geq 70\%$ & Befriedigend \\
    2,3 & $\geq 69$ & $\geq 63\%$ & Befriedigend \\
    2,7 & $\geq 60$ & $\geq 55\%$ & Ausreichend \\
    3,0 & $\geq 55$ & $\geq 50\%$ & Ausreichend \\
    4,0 & $< 55$ & $< 50\%$ & Nicht bestanden \\
    \hline
\end{tabular}

\subsection*{Wichtige Korrektor-Hinweise}

\begin{enumerate}[1.]
    \item \textbf{Teilpunkte großzügig vergeben:}
    \begin{itemize}
        \item Auch wenn Lösung nicht 100\% korrekt ist
        \item Wenn Verständnis der Konzepte erkennbar ist
    \end{itemize}
    
    \item \textbf{Rechenfehler:}
    \begin{itemize}
        \item Wenn Rechenweg richtig ist: volle Punkte trotz Rechenfehler
        \item Nur Punkte abziehen wenn Ansatz falsch ist
    \end{itemize}
    
    \item \textbf{Alternative Lösungen:}
    \begin{itemize}
        \item Akzeptiere alternative korrekte Antworten
        \item z.B. andere sichere Verschlüsselungsmodi (OFB, GCM statt CBC)
        \item Andere Begründungen, solange sachlich korrekt
    \end{itemize}
    
    \item \textbf{Bonuspunkte:}
    \begin{itemize}
        \item Für besonders tiefe Einsichten
        \item Für praktische Beispiele
        \item Für Referenzen zu aktuellen Sicherheitsvorfällen
    \end{itemize}
\end{enumerate}

\end{document}
